\section{Bohr sets and partitions}
\label{sec:partition}

\subsection{Monochromatic configurations}
We make some preparations before the proof of \cref{th:counting-partition}. In this section, we only need $G$ to be a commutative semigroup with neutral element. Fix $k+1$ commuting (semigroup)  homomorphisms $\psi, \phi_1, \ldots, \phi_k: G \rightarrow G$ with $\psi \neq 0$. 
We write 
\[
\Phi_m = \{ \psi^{i_0} \circ \phi_1^{i_1} \circ \cdots \circ \phi_k^{i_k} : 0 \leq i_0, i_1, \ldots, i_k \leq m \} \cup \{ 0 \}
\] (where $\phi^i$ is the $i$-th composition of $\phi$).

For \textit{formal variables} $x_1, \ldots, x_n$, we write
\[
S_m(x_1,\ldots, x_n) = \left\{ \sum_{i=1}^n \xi_i (x_i) : \xi_i \in \Phi_m \right\} 
\]
and we refer to $S_m(x_1,\ldots, x_n)$ as the $S_{m,n}$-\textit{set} with \textit{generators} $x_1, \ldots, x_n$. There is a canonical bijection between $S_m(x_1, \ldots, x_n)$ and $\Phi_m^n$ defined by
\[
    \sum_{i=1}^n \xi_i(x_i) \mapsto (\xi_1, \ldots, \xi_n).
\]
 For an element $x = \sum_{i=1}^n \xi_i (x_i) \in S_m(x_1,\ldots, x_n)$, by the \textit{support} of $x$ we mean the set $\{ i \in [n]: \xi_i \neq 0\}$.
The goal of this section is to prove the following:

\begin{theorem} \label{th:mpc-hom} 
For any $r > 0$, there exist $n$ and $m$ such that under any $r$-coloring of $S_m(x_1,\ldots, x_n)$, there is a monochromatic configuration
\[
\{ \psi(y), x, x + \phi_1(y), \ldots, x+ \phi_k(y) \},
\]
where $x, y$ have nonempty and disjoint supports.
%\Hnote{This will help us when we integrate and change variables}.
\end{theorem}
The fact that the supports of $x$ and $y$ are nonempty and disjoint will be crucial in our applications (\cref{th:counting-partition} and \cref{prop:counting-brauer}). \cref{th:mpc-hom} follows from \cref{prop:mpc} below whose proof requires the multidimensional Hales-Jewett theorem (for a reference, see \cite[Theorem 7, p.40]{ramsey-book}). We recall the theorem here for reader's convenience.

%Before stating the theorem, we need to \cref{th:mpc-hom} follows from \cref{prop:mpc}. The main ingredient in the proof of \ref{prop:mpc} multidimensional Hales-Jewett theorem, which we now recall (see e.g. \cite[Theorem 7, p.40]{ramsey-book}). 

The set $[t]^N = \{(x_1, \ldots, x_N): x_i \in [t] \}$ is called a \textit{cube} of dimension $N$ over $t$ elements. Let $[N] = A_0 \cup A_1 \cup \cdots \cup A_m$ be any disjoint partition of $[N]$, where $A_i \neq \varnothing$ for $i \neq 0$ ($A_0$ may be empty), and $f: A_0 \rightarrow [t]$ be any map. Define a map $g: [t]^m \rightarrow [t]^N$ by assigning to each $(y_1, \ldots, y_m) \in [t]^m$ the element $(x_1, \ldots, x_N) \in [t]^N$, where
\begin{equation}
x_i = 
\begin{cases}
f(i), \quad &\textup{if}\ i \in A_0 \\
y_j, \quad &\textup{if}\ i \in A_j \text{ for }j \in [m].
\end{cases}
\end{equation}
A \textit{combinatorial space} of dimension $m$ is the image of $g$ for some choice of $A_0, A_1, \ldots, A_m$ and $f$. We can now state:

\begin{theorem}[Multidimensional Hales-Jewett] For any $r, t, m$, there exists a number $N=HJ(t, m;r)$ such that whenever $[t]^N$ is $r$-colored, there must be a monochromatic combinatorial space of dimension $m$.
\end{theorem}

Using this, we can prove the following proposition. (Recall that $k$ is fixed from the beginning of this section.)

\begin{proposition} \label{prop:mpc}
 For any $r>0$ and $\ell >0$, there exist $n=n(k,\ell, r)$ and $m=m(k,\ell, r)$ such that under any $r$-coloring of $S_{m}(x_1,\ldots, x_n)$, there are elements 
$y_1, \ldots, y_\ell \in S_m(x_1,\ldots, x_n)$ with nonempty and disjoint supports, such that for each $i \in [\ell]$, the elements
\[
\psi(y_i) + \sum_{1 \leq j \leq i-1} \xi_j (y_j) \qquad \textup{ where } \xi_j \in \{ 0, \psi, \phi_1, \ldots, \phi_k \}
\]
have the same color (i.e. their color depends only on $i$).
\end{proposition}

\begin{proof}
The number of colors $r$ will be fixed throughout. We will proceed by induction on $\ell$. When $\ell=1$ the statement is obvious. Suppose the statement is true for $\ell$, we will prove it is true for $\ell+1$. 

Write $n' = n(k,\ell, r), m'=m(k,\ell, r)$. We define $m = m(k,\ell + 1, r):=|\Phi_{m'+1}| +1$, $N:=HJ( |\Phi_{m'+1}|, n';r)$ and $n = n(k, \ell+1,r): = 1 + N$.

Consider an arbitrary $r$-coloring of $S_{m}(x_1,\ldots, x_n)$. An $r$-coloring of $S_{m}(x_1,\ldots, x_n)$ induces an $r$-coloring of $\Phi_{m'+1}^{N}$ by assigning to 
$(a_1, \ldots, a_{N}) \in (\Phi_{m'+1})^{N}$ the color of 
\[
\psi (x_n) + \sum_{i=1}^{N} \psi \circ a_i (x_i).
\]
Since $N=HJ( |\Phi_{m'+1}|, n';r)$, there are a disjoint partition
\[
[N] = A_0 \cup A_1 \cup \cdots \cup A_{n'}, \quad A_i \neq \varnothing, \forall i \neq 0
\]
and functions $f_i \in \Phi_{m' + 1}$ for $i \in A_0$ such that when $\zeta_1, \zeta_2, \ldots, \zeta_{n'}$ range over $\Phi_{m'+1}$, all the elements 
\[
\psi(x_n) + \sum_{i=1}^{N} \psi \circ a_i (x_i),
\] 
with
\begin{equation*}
a_i = 
\begin{cases}
f_i, \quad &\textup{if}\ i \in A_0 \\
\zeta_j, \quad &\textup{if}\ i \in A_j \text{ for } 1 \leq j \leq n'
\end{cases}
\end{equation*}
have the same color.

Write $z_j = \sum_{i \in A_j} \psi(x_i)$ for $1 \leq j \leq n'$, and $z_{n'+1} = x_{n} + \sum_{i \in A_0} f_i(x_i)$. Then all the $z_j$ have nonempty and disjoint supports, and all elements of the form
\[
\psi(z_{n'+1}) + \sum_{j=1}^{n'} \zeta_j (z_j), \qquad \zeta_j \in \Phi_{m'+1},
\]
have the same color.

%Clearly,
%\[
%S_{k+1}(x_1,\ldots, x_n) \supset S_k (z_1, \ldots, z_n').
%\]
By the inductive hypothesis, there exists a sequence $y_1, \ldots, y_\ell \in S_{m'} (z_1, \ldots, z_{n'})$ having nonempty and disjoint supports such that for each $i=1, \ldots, \ell$, the elements
\[
\psi(y_i) + \sum_{1 \leq j \leq i-1} \xi_j (y_j) \qquad \textup{ where } \xi_j \in \{ 0, \psi, \phi_1, \ldots, \phi_k \}
\]
have the same color. We now set $y_{\ell+1} = z_{n'+1}$. Clearly the elements 
\[
\psi(y_{\ell+1}) + \sum_{1 \leq j \leq \ell} \xi_j (y_j) \qquad \textup{ where } \xi_j \in \{ 0, \psi, \phi_1, \ldots, \phi_k \}
\]
are of the form
\[
\psi(z_{n'+1}) + \sum_{j=1}^{n'} \zeta_j (z_j), \qquad \zeta_j \in \Phi_{m'+1},
\]
and so they have the same color. Thus  Proposition \ref{prop:mpc} is proved.
\end{proof}

\begin{proof}[Proof of \cref{th:mpc-hom}]
Applying Proposition \ref{prop:mpc} with $\ell=r+1$, we can find a sequence $y_1, \ldots, y_{r+1}$ satisfying the conclusion of that proposition. Let $c(i)$ be the color of
\[
\psi(y_i) + \sum_{1 \leq j \leq i-1} \xi_j (y_j) \qquad \textup{ where } \xi_j \in \{ 0, \psi, \phi_1, \ldots, \phi_k \}.
\]
Then there exist $1 \leq u < v \leq r+1$ such that $c(u) = c(v)$. Hence the elements
\[
\psi(y_u), \psi(y_v), \psi(y_v) + \phi_1(y_u), \ldots, \psi(y_v) + \phi_k(y_u), 
\]
have the same color, and we are done (with $x=\psi(y_v), y=y_u$).
\end{proof}


\subsection{Proofs of \texorpdfstring{\cref{th:counting-partition}}{Theorem 1.11} and \texorpdfstring{\cref{th:main-partition}}{Theorem 1.5}}

Using Theorem \ref{th:mpc-hom} we can now prove \cref{th:counting-partition}, which we recall for convenience:

\begin{theorem*} 
Suppose $\psi, \phi_1, \ldots, \phi_k: G \to G$ are continuous homomorphisms satisfying:
\begin{enumerate}
\item $\psi, \phi_1, \ldots, \phi_k$ are commuting, and

\item $\psi(G), \phi_1(G), \ldots, \phi_k(G)$ have finite index in $G$.
\end{enumerate}
Suppose $G = \bigcup_{i=1}^r A_i$ is a partition of $G$ into measurable sets. Then 
\begin{equation*} 
\sum_{i=1}^r \iint_{G^2} 1_{A_i}(\psi(t)) 1_{A_i}(x) 1_{A_i}(x + \phi_1(t)) \cdots 1_{A_i}(x + \phi_k(t)) \, d \mu(x) d\mu(t) \geq c_3
\end{equation*} 
for some positive constant $c_3$ depending only on $r, k$ and the aforementioned indexes.
\end{theorem*}

\begin{proof}Consider the set $S_m(x_1, \ldots, x_n)$ given by Theorem \ref{th:mpc-hom}, where we now let $x_1, \ldots, x_n$ vary over $G$. Note that for any non-zero $\phi \in \Phi_m$, we have $[G:\phi(G)] < \infty$ by \cref{lem:composition}. Let $R=R(x_1, \ldots, x_n)$ be the set of all pairs $(z, t)$ where $z, t \in S_m(x_1, \ldots, x_n)$ have nonempty and disjoint supports.

Suppose $G = \bigcup_{i=1}^r A_i$. For $i \in [r]$, we define
\[
T_i : = \iint_{G^2} 1_{A_i}( \psi(y) )  1_{A_i}(x) 1_{A_i}(x+ \phi_1(y)) \cdots 1_{A_i}(x+ \phi_k(y)) \, d\mu(x) d\mu(y). 
\]
Let $(z,t) \in R$ be arbitrary, and suppose
\[
z = \sum_{u \in U} \zeta_u(x_u) \qquad \textup{and} \qquad t = \sum_{v \in V} \xi_v (x_v)
\]
where $U, V \subset [n]$ are nonempty and disjoint and $\zeta_u, \xi_v \in \Phi_m \setminus \{ 0\}$. We have
\begin{eqnarray*}
& & \int_{G^n}  1_{A_i}( \psi(t)) 1_{A_i}(z) 
1_{A_i}(z+ \phi_1(t)) \cdots 1_{A_i}(z+ \phi_k(t)) \, d\mu(x_1) \cdots d\mu(x_n) \\
& = & \int_{G^n} 1_{A_i} \left( \sum_{v \in V} \psi( \xi_v (x_v) ) \right) 
1_{A_i} \left( \sum_{u \in U} \zeta_u(x_u)\right)
1_{A_i} \left( \sum_{u \in U} \zeta_u(x_u)  + \phi_1 \left( \sum_{v \in V} \xi_v (x_v) \right) \right) \cdots \\
& & \qquad 1_{A_i} \left( \sum_{u \in U} \zeta_u(x_u)  + \phi_k \left( \sum_{v \in V} \xi_v (x_v) \right)  \right) \, d\mu(x_1) \cdots d \mu(x_n) \\
& \ll & \int_{G^n} 1_{A_i} \left( \sum_{v \in V} \psi( x_v ) \right) 
1_{A_i} \left( \sum_{u \in U} x_u \right)
1_{A_i} \left( \sum_{u \in U} x_u  + \phi_1 \left( \sum_{v \in V} x_v \right) \right) \cdots \\
& & \qquad 1_{A_i} \left( \sum_{u \in U} x_u  + \phi_k \left( \sum_{v \in V} x_v \right)  \right) \, d\mu(x_1) \cdots d \mu(x_n) \\
& = & \iint_{G^2} 1_{A_i}( \psi(y) ) 1_{A_i}(x) 1_{A_i}(x  + \phi_1(y) ) \cdots 1_{A_i}(x  + \phi_k(y) ) \, d\mu(x) \, d\mu(y) \\
&=& T_i, 
\end{eqnarray*}
by $|U|+|V|$ applications of \cref{lem:finite-index}. 


Fixing any $(x_1, \ldots, x_n) \in G^n$, the set $S_m(x_1, \ldots, x_n)$ becomes a subset of $G$ and so the coloring $G = \bigcup_{i=1}^r A_i$ naturally induces a coloring of $S_m(x_1, \ldots, x_n)$. By \cref{th:mpc-hom}, there exists $i \in [r]$ and $(z, t) \in R$ such that
\[
    \psi(t), z, z+ \phi_1(t), \cdots, z+ \phi_k(t) \in A_i.
\]
Thus, for every $(x_1, \ldots, x_n) \in G^n$,
\[
    \sum_{i=1}^r \sum_{(z, t) \in R } 1_{A_i}( \psi(t)) 1_{A_i}(z) 1_{A_i}(z+ \phi_1(t)) \cdots 1_{A_i}(z+ \phi_k(t)) \geq 1.
\]
It follows that
\begin{eqnarray*}
1 & \leq & \int_{G^n} \sum_{i=1}^r \sum_{(z, t) \in R } 1_{A_i}( \psi(t)) 1_{A_i}(z) 1_{A_i}(z+ \phi_1(t)) \cdots 1_{A_i}(z+ \phi_k(t)) \, d\mu(x_1) \cdots d\mu(x_n) \\
&\leq& \sum_{i=1}^r  \int_{G^n} \sum_{(z, t) \in R } 1_{A_i}( \psi(t)) 1_{A_i}(z) 1_{A_i}(z+ \phi_1(t)) \cdots 1_{A_i}(z+ \phi_k(t)) \, d\mu(x_1) \cdots d\mu(x_n)  \\
& \ll & \sum_{i=1}^r T_i,   
\end{eqnarray*}
thus finishing the proof.
\end{proof}

% \textcolor{red}{Rewrite this sentence.}
To prove \cref{th:main-partition}, we will need the following proposition. 
% need only the case $k=1$ of \cref{th:counting-partition}.  %\textcolor{red}{It's a bit awkward here when we prove for $f_i$ while later we only use $1_{A_i}$.}
With an eye to potential applications, we state and prove a slightly stronger version than what is needed.

%The following proposition is slightly stronger than what we need. However, since its proof is not harder, we give it here for future applications.


%We will record the following fact for future applications.

\begin{proposition} \label{prop:partition-bohr-stronger} 
Let $\phi, \psi: G \to G$ be commuting continuous homomorphisms with images of finite index. Suppose $f_1, \ldots, f_r: G \to [0,1]$ are measurable functions such that $\sum_{i=1}^r f_i \geq 1$ pointwise. For $w \in G$, define
\[
  R_i(w) =  \iint_{G^2}  f_i(\psi(y)) f_i(x+w) f_i(x + \phi(y)) \ d\mu(x) d\mu(y). 
\]
Then there are $c, k, \eta >0$ depending only on $r$ and the indexes above such that for some $i \in [r]$, the set
 $\{ w \in G: R_i(w) > c\}$ contains a Bohr-$(k,\eta)$ set. 
\end{proposition}
\begin{proof}
%The proof is the same as that of \cref{prop:partition-bohr-stronger}, using \cref{lem:counting2} instead of \cref{lem:bourgain_lemma_2}. \Hnote{to fill out}

%\textcolor{red}{Fix the first part a bit}.
For $i \in [r]$, let $A_i = \{x \in G: f_i(x) \geq 1/r\}$. Since $\sum_{i=1}^r f_i \geq 1$ pointwise, $G = \bigcup_{i=1}^r A_i$. In light of \cref{th:counting-partition}, there exists a constant $c$ depending only on $r$ and the indexes and an $i \in [r]$ such that 
\[
    \iint_{G^2} 1_{A_i} (\psi(y)) 1_{A_i}(x) 1_{A_i}(x + \phi(y)) \ d\mu(x) d\mu(y) > c.
\]
It then follows that
\[
    R_i(0) \geq \frac{c}{r^3}.  
\]
On the other hand, by \cref{lem:counting2}, for every $w \in G$,
\begin{equation*}
\label{eq:partition_R_i}
    |R_i(w) - R_i(0)| \ll \lVert \widehat{f_i} - \widehat{f_{i,w}} \rVert_{\infty},
\end{equation*}
where the implicit constant depends only on the indexes of $\phi(G)$ and $\psi(G)$ in $G$. 
%In light of \cref{th:counting-partition}, there exists $i_0 \in [r]$ such that $R_{i_0}(0) > c$ for some constant $c>0$, where $c$ depends only on $r$ and the aforementioned indices. 
Hence, there exists a constant $c'$ such that $R_{i}(w) \geq \frac{c}{2r^3}$ if
\begin{equation*}
% \label{eq:br_f_w_needed}
    \lVert \widehat{f_{i}} - \widehat{f_{i,w}} \rVert_{\infty} < c'.
\end{equation*}
By \cref{lem:translate}, the set of such $w$ contains a Bohr-$(k,\eta)$ set, where $k$ and $\eta$ depend only on $c'$.
\end{proof}
 
\cref{th:main-partition} is now a special case of the next theorem with $\psi_1 = \phi_2$ and $\psi_2 = \phi_1$. 

\begin{theorem}\label{th:main-partition-true}
Let $G=\bigcup_{i=1}^r A_i$ be a partition into measurable sets. Let $\phi_1, \phi_2, \psi_1, \psi_2: G \rightarrow G$ be continuous homomorphisms satisfying the following:
\begin{enumerate}
\item $\phi_2 \circ \psi_2 = \phi_1 \circ \psi_1$,
\item $\psi_1 \circ \psi_2 = \psi_2 \circ \psi_1$,
\item $\phi_1(G), \psi_1(G), \psi_2(G)$ have finite index in $G$.
\end{enumerate}
Then for some $1 \leq i \leq r$, the set $\phi_1(A_i) - \phi_1(A_i) + \phi_2(A_i)$ contains a Bohr-$(k, \eta)$ set, where $k$ and $\eta$ depend only on $r$ and the indexes of $\phi_1(G), \psi_1(G), \psi_2(G)$ in $G$.
\end{theorem}

\begin{proof}
Suppose $G = \bigcup_{i=1}^r A_i$. We apply \cref{prop:partition-bohr-stronger} with $f_i = 1_{A_i}$ and $(\psi_1, \psi_2)$ in place of $(\psi, \phi)$. Then for some $i$, the set $\{ w \in G: R_i(w) > c\}$ contains a Bohr-$(k,\eta)$ set $B$. This means that for $w \in B$, there exist $x, y \in G$ such that
\[
x+w, \psi_2(y), x + \psi_1(y) \in A_i.
\]
Since
\[
\phi_1(x+w) + \phi_2( \psi_2(y)) - \phi_1( x + \psi_1(y)) = \phi_1(w),
\]
we conclude that $\phi_1(B) \subset \phi_1(A_i) + \phi_2(A_i) - \phi_1(A_i)$. Our theorem now follows from \cref{lem:bohr-homo2}.
\end{proof}


\begin{remark} \label{rk:auto2}
Here we explain why the commutativity condition of $\phi_1$ and $\phi_2$ in \cref{th:main-partition} is not necessary if $\phi_1$ or $\phi_2$ is an automorphism (see \cref{remark:first_one}). If $\phi_1$ is an automorphism, in \cref{th:main-partition-true}, we take $\psi_1 = \phi_1^{-1} \circ \phi_2$ and $\psi_2 = \textup{Id}$, the identity homomorphism. Then the first two conditions of \cref{th:main-partition-true} are satisfied. As for the third condition, we have $\psi_1(G) = \phi_1^{-1} \circ \phi_2(G),$ which has finite index in $G$ by \cref{lem:composition}. A similar argument applies to the case $\phi_2$ is an automorphism.
\end{remark}
